% =====================================================================
% DiVa -- Die Mathematik der Messung
% Ein in sich geschlossener Bericht (deutsch), kompilierbar mit pdflatex.
% Alle Zahlenwerte sind direkt aus results/*.json reproduzierbar.
% =====================================================================
\documentclass[11pt,a4paper]{article}
\usepackage[margin=2.4cm]{geometry}
\usepackage{amsmath,amssymb,amsthm}
\usepackage[T1]{fontenc}
\usepackage[utf8]{inputenc}
\usepackage[ngerman]{babel}
\usepackage{booktabs}
\usepackage{microtype}

\newtheorem{satz}{Satz}
\newtheorem{lemma}[satz]{Lemma}
\newtheorem{korollar}[satz]{Korollar}
\newtheorem{folgerung}[satz]{Folgerung}
\newtheorem{definition}[satz]{Definition}
\newtheorem{befund}{Empirischer Befund}
\newtheorem{vermutung}{Offenes Problem}
\theoremstyle{remark}
\newtheorem{bemerkung}{Bemerkung}

\newcommand{\Gr}{\mathrm{Gr}}
\newcommand{\E}{\mathbb{E}}
\newcommand{\R}{\mathbb{R}}
\newcommand{\ind}[1]{\mathbf{1}_{#1}}

\title{DiVa: Die Mathematik der gemessenen kritischen Fläche\\
\large Swarm-Synchronisation auf der Graßmann'schen Mannigfaltigkeit\\
unter Anker-Interpolation und Rauschen}
\author{Autonome numerische Sitzung (GPU-Verifikationsprogramm)}
\date{15.\,November 2026}

\begin{document}
\maketitle

\begin{abstract}
Wir betrachten ein Schwarm-Modell auf der Graßmann-Mannigfaltigkeit
$\Gr(k,d)$ mit Mittelfeld-Kopplung $\eta$, tangencialem Gauß-Rauschen der
Stärke $\sigma$ und einem stochastischen Anker-Term
$\lambda_W\,p_{\mathrm{fix}}$. Die ursprüngliche Konjektur postuliert eine
kritische Fläche
$\eta_c(\sigma,\lambda_W,p_{\mathrm{fix}})=\eta_c(0)-a_\sigma\sigma^2-a_W\lambda_W^2p_{\mathrm{fix}}$.
Die numerische Verifikation bestätigt die \emph{quadratische Form} in
$\sigma^2$, widerlegt aber beide Koeffizienten-Gesetze und liefert als
Ersatz drei einfache, robuste Gesetze:
(i)~$\eta_{\mathrm{lo}}(\sigma)=a_\sigma\sigma^2$ mit
$a_\sigma(k,d)=3d+k$ (12 Konfigurationen, $\pm5\,\%$, float64-bestätigt),
(ii)~ein geometrieunabhängiges oberes Fensterende
$\eta_{\mathrm{hi}}\approx2{,}04$ für $k\ge2$,
(iii)~eine universelle Hoeffding-Signatur $\zeta_1^\ast\approx0{,}23$--$0{,}25$
am Rand. Die Marginal-Entartung der geordneten Phase wird operationell als
Lyapunov-Plateau direkt nachgewiesen ($d\ge4$) und von Ankern angehoben.
Dieses Dokument stellt Modell, Definitionen, Sätze und alle Messwerte
so zusammen, dass sie in einem Durchgang lesbar sind.
\end{abstract}

\tableofcontents

% =====================================================================
\section{Das Modell}

\begin{definition}[Konfigurationsraum]\label{def:gr}
Die Graßmann-Mannigfaltigkeit $\Gr(k,d)$ ist die Menge der
$k$-dimensionalen Unterräume $V\subset\R^d$. Wir repräsentieren
$V\in\Gr(k,d)$ durch Matrizen $Q\in\R^{d\times k}$ mit
$Q^\top Q=I_k$; zwei Repräsentatoren beschreiben denselben Punkt genau
dann, wenn sie sich um eine rechtsseitige orthogonale Matrix
unterscheiden ($Q\sim QO$, $O\in O(k)$). Der Tangentialraum in $Q$ ist
\[
T_Q\Gr(k,d)\;=\;\{\,\xi\in\R^{d\times k}\;:\;Q^\top\xi=0\,\}
\;=\;(I-QQ^\top)\,\R^{d\times k}.
\]
\end{definition}

\begin{definition}[Interpolierter-Anker-Schwarm]\label{def:modell}
Gegeben $N$ Agenten $Q_1,\dots,Q_N\in\Gr(k,d)$, den Mittelfeld-Projektor
\[
\bar S \;=\;\frac1N\sum_{j=1}^N Q_j Q_j^\top \;\in\;\R^{d\times d},
\]
Parameter $\eta,\lambda_W,\sigma\ge0$, $p_{\mathrm{fix}}\in[0,1]$, einen
\emph{gelöschten} (quenched) Anker $W_{\mathrm{fix}}\sim\mathrm{Haar}$ auf
$\Gr(k,d)$ sowie für jeden Zeitschritt unabhängige Zufallsgrößen
$W_i^{(t)}$ mit
\[
W_i^{(t)}=\begin{cases}W_{\mathrm{fix}} & \text{mit Wahrsch.~}p_{\mathrm{fix}},\\
\mathrm{Haar} & \text{sonst,}\end{cases}
\qquad
G_i^{(t)}\sim\mathcal N(0,I_{d\times k}).
\]
Ein Updateschritt ist
\begin{equation}\label{eq:dyn}
Q_i^{(t+1)}=\mathrm{orth}\Big(
Q_i^{(t)}
+\eta\,(I-Q_iQ_i^{\top})\,\bar S\,Q_i^{(t)}
+\lambda_W\,(I-Q_iQ_i^{\top})\,W_i^{(t)}W_i^{(t)\top}Q_i^{(t)}
+\sigma\,\underbrace{(I-Q_iQ_i^{\top})\,G_i^{(t)}}_{=:Z_i^{(t)}}
\Big),
\end{equation}
wobei $\mathrm{orth}(X)$ irgendein orthonormales Repräsentator-Polynom der
Spalten von $X$ ist (wir benutzen Gram--Schmidt).
\end{definition}

\begin{bemerkung}[Skalierung des Rauschens]
Für $Z=(I-QQ^\top)G$ gilt $\E[\mathrm{vec}(Z)\mathrm{vec}(Z)^\top]
=(I-QQ^\top)\otimes I_k$: jede der $k(d-k)$ Tangentialrichtungen erhält
Varianz~$1$, mithin ist $\sigma Z$ ein tangentialer Gauß-Schub der
Stärke $\sigma$ pro Freiheitsgrad; ferner
$\E\|Z\|_F^2=k(d-k)$.
\end{bemerkung}

\begin{lemma}[Wohldefiniertheit auf dem Quotienten]\label{lem:quot}
Die Abbildung \eqref{eq:dyn} hängt nur von den Projektoren
$P_i=Q_iQ_i^\top$ ab und vertauscht mit der rechtsseitigen
$O(k)$-Wirkung; sie induziert also eine wohldefinierte Abbildung
$F\colon\Gr(k,d)^N\to\Gr(k,d)^N$.
\end{lemma}
\begin{proof}
Für $Q' = QO$ gilt $Q'Q'^\top = QQ^\top$, also
$(I-Q'Q'^\top)=(I-QQ^\top)$, und jeder Driftterm transformiert kovariant,
$(I-QQ^\top)\bar S\,(QO) = \big[(I-QQ^\top)\bar S\,Q\big]\,O$ (entsprechend
Anker und Rauschen). Der gesamte Inkrement hat also die Form
$\Delta(QO)=\Delta(Q)\,O$. Wegen $\mathrm{span}(\mathrm{orth}(X))=\mathrm{span}(X)$
folgt
$\mathrm{span}\big(F(QO)\big)=\mathrm{span}\big(QO+\Delta(Q)O\big)
=\mathrm{span}\big(Q+\Delta(Q)\big)\cdot O$, d.\,h.\ die Abbildung auf
Subräumen ist wohldefiniert und $O(k)$-äquivariant.
\end{proof}

Alle folgenden Aussagen betreffen den irreduziblen Stationärzustand
(NESS) der Kette $Q^{(t+1)}=F(Q^{(t)})$ auf $\Gr(k,d)^N$.

% =====================================================================
\section{Beobachtungsgrößen}

\begin{definition}[Geodätischer Abstand]\label{def:geo}
Für $V,W\in\Gr(k,d)$ seien $\theta_1,\dots,\theta_k\in[0,\pi/2]$ die
Hauptwinkel, definiert durch
$\cos\theta_m = s_m(V^{\top}W)$ mit den Singulärwerten
$s_m$ von $Q_V^\top Q_W$. Der geodätische Abstand (für die Standard-
Quotientenmetrik) ist
\[
d_G(V,W)\;=\;\Big(\sum_{m=1}^k \theta_m^2\Big)^{1/2}.
\]
\end{definition}

\begin{definition}[Ordnungsparameter]
\begin{align}
\Delta_T &\;=\;\frac{2}{N(N-1)}\sum_{i<j} d_G(V_i,V_j)
&&\text{(U-Statistik der Ordnung 2, Kern $h=d_G$)},\\
R &\;=\;\lambda_{\max}(\bar S)
&&\text{(Mean-Feld-Parameter: $\approx k/d$ ungeordnet, $\approx1$ geordnet)}.
\end{align}
\end{definition}

\begin{definition}[Kritische Kanten, operationell]
Für feste $(k,d,N,\sigma,\lambda_W,p_{\mathrm{fix}})$ sei
\[
\Delta(\eta)\;=\;\text{Zeitmittel von $\Delta_T$ im NESS bei Kopplung }\eta .
\]
$\Delta(\eta)$ hat (für $\sigma>0$) die Gestalt
,,Hoch--Tief--Hoch'': Plateau $\Delta_{\mathrm{dis}}$ bei kleinem $\eta$,
Einbruch auf $\approx0$ im Fenster, Wiederanstieg bei großem $\eta$.
Wir definieren
\[
\eta_{\mathrm{lo}}=\inf\{\eta:\Delta(\eta)<\tfrac12\Delta_{\mathrm{dis}}\},
\qquad
\eta_{\mathrm{hi}}=\sup\{\eta:\Delta(\eta)<\tfrac12\Delta_{\mathrm{dis}}\},
\]
mit logarithmischer Interpolation zwischen Gitterpunkten
(Protokoll: $T\ge6000$ Schritte pro Punkt, $\ge8$ Zeitmittel,
Replikate $P\ge8$; T-Konvergenz $\le1\,\%$ zwischen $T=4000$ und
$T=12000$; Replikat-Streuung $\sigma\le2\,\%$).
\end{definition}

\subsection{Die erste Hoeffding-Zerlegung und $\zeta_1$}

\begin{satz}[Hoeffding-Zerlegung]\label{thm:hoeffding}
Es seien $X_1,\dots,X_N$ i.i.d.\ mit Werten in einem polnischen Raum und
$h\ge0$ ein symmetrischer Kern mit $\E h^2<\infty$. Mit
\[
\theta=\E h(X_1,X_2),\qquad
g_1(x)=\E\big[h(x,X_1)\big]-\theta,\qquad
\sigma_h^2=\mathrm{Var}\,h(X_1,X_2),
\]
gilt für die U-Statistik
\[
U_N-\theta \;=\;\frac{2}{N}\sum_{i=1}^N g_1(X_i)\;+\;R_N,
\qquad \E R_N=0,\quad R_N=O_p(N^{-1}),
\]
und $\mathrm{Var}\big(\tfrac2N\sum_i g_1(X_i)\big)=\dfrac{4\,\zeta_1\,\sigma_h^2}{N}$
mit dem \emph{ersten Hoeffding-Koeffizienten}
\[
\zeta_1\;=\;\frac{\mathrm{Var}\,g_1(X_1)}{\sigma_h^2}\;\in\;[0,1].
\]
\end{satz}

\begin{korollar}[Operationaler Schätzer für $\zeta_1$]\label{kor:zeta}
Definiere Zeilenmittel $r_i=\frac{1}{N-1}\sum_{j\ne i}h(X_i,X_j)$. Dann
\[
r_i \;=\;\theta+g_1(X_i)+\varepsilon_i,\qquad
\E[\varepsilon_i\mid X_i]=0,\qquad
\mathrm{Var}(\varepsilon_i\mid X_i)=\frac{\sigma_h^2(X_i)}{N-1},
\]
also
\[
\mathrm{Var}_i(r_i)\;=\;\zeta_1\,\sigma_h^2\;+\;\frac{\sigma_h^2}{N-1}\;+\;o(N^{-1}).
\]
Der konsistente Schätzer ist folglich
\begin{equation}\label{eq:zetaest}
\widehat{\zeta_1}\;=\;\frac{\max\!\Big(0,\ \mathrm{Var}_i(r_i)-\dfrac{\widehat\sigma_h^2}{N-1}\Big)}{\widehat\sigma_h^2}.
\end{equation}
\end{korollar}
\begin{proof}
Bedingung auf $X_i$: $\E[r_i\mid X_i]=g_1(X_i)+\theta$ und die bedingte
Varianz des Stichprobenmittels der $h(X_i,X_j)$, $j\ne i$, ist
$\sigma_h^2(X_i)/(N-1)$; der Kreuzterm verschwindet, weil
$\E[\varepsilon_i\mid X_i]=0$. Mittelung über $X_i$ liefert die Formel mit
$\E[\sigma_h^2(X_i)]=\sigma_h^2$.
\end{proof}

\begin{bemerkung}
Gleichung \eqref{eq:zetaest} ist genau die in der Simulation benutzte
Größe: $\zeta_1=0$ bedeutet, dass die Agentenbezüge $r_i$ bis auf
Stichprobenrauschen konstant sind (,,flache erste Hoeffding-Richtung''),
$\zeta_1>0$ bedeutet sichtbare Heterogenität (z.\,B.\ Clusterstruktur).
\end{bemerkung}

% =====================================================================
\section{Die Phasenstruktur: das Dichotomie-Fenster}

\begin{befund}[Dichotomie-Fenster]\label{bef:fenster}
Für $\lambda_W=0$ existiert zu jedem $\sigma\ge0$ ein Intervall
$[\eta_{\mathrm{lo}}(\sigma),\eta_{\mathrm{hi}}]$ mit:
\begin{enumerate}
\item[(i)] $\eta<\eta_{\mathrm{lo}}$: \emph{ungeordnete NESS},
$\Delta\approx\Delta_{\mathrm{dis}}(\sigma)>0$
(z.\,B.\ $\Delta_{\mathrm{dis}}=1{,}59$ für $(k,d)=(2,8)$, $N=64$);
\item[(ii)] $\eta_{\mathrm{lo}}<\eta<\eta_{\mathrm{hi}}$: \emph{geordnete Phase},
$\Delta\le3\cdot10^{-4}$ (bei $N=64$), $R=1$;
\item[(iii)] $\eta>\eta_{\mathrm{hi}}$: \emph{Re-Desynchronisation},
$\Delta=O(1)$ mit chaotischer Fluktuation.
\end{enumerate}
Ohne Rauschen ($\sigma=0$) und ohne Anker ist der exakt synchronisierte
Zustand für \emph{alle} $\eta$ im Fenster ein invarianter Zustand
(dort verschwindet der Drift); die Kanten sind dann
stabilitäts-induzierte Objekte, kein Plateau-Übergang.
\end{befund}

% =====================================================================
\section{Das $\sigma^2$-Gesetz und $a_\sigma=3d+k$}

\begin{befund}[$\sigma^2$-Gesetz]\label{bef:sigma}
Für $\lambda_W=0$ gilt im Bereich $0{,}02\le\sigma\le0{,}2$:
\[
\eta_{\mathrm{lo}}(\sigma;k,d)\;=\;a_\sigma(k,d)\,\sigma^2,\qquad
\text{Achsenabschnitt } \eta_0 = 0{,}000\pm0{,}001 .
\]
Beispiel $(2,8)$: $\eta_{\mathrm{lo}}/\sigma^2 = 25{,}1,\,25{,}2,\,25{,}5,\,25{,}6,\,26{,}2$
für $\sigma=0{,}02,\,0{,}03,\,0{,}05,\,0{,}07,\,0{,}1$ --- konstant auf
$\pm2\,\%$. Erst wenn $\eta_{\mathrm{lo}}$ in die Nähe von
$\eta_{\mathrm{hi}}/2$ rückt (Fenster-Schluss-Regime), driftet das Verhältnis
nach oben (28{,}5 bzw.\ 35{,}7) --- kein Bruch des Gesetzes, sondern
Endeneffekt. Das Gesetz ist $N$-unabhängig:
$a_\sigma=25{,}12,\,25{,}47,\,25{,}11$ für $N=32,64,128$.
\end{befund}

\begin{befund}[Koeffizientengesetz]\label{bef:3dk}
Über 12 Konfigurationen (je zwei unabhängige $\sigma$-Werte, float64 wo
angegeben):
\begin{center}
\begin{tabular}{lrrr}
\toprule
$(k,d)$ & $a_\sigma$ gemessen & $3d+k$ & Verhältnis\\
\midrule
(1,2) & $\mathbf{7{,}00\pm0{,}14}$ (float64) & 7 & 1{,}00\\
(1,3) & $9{,}8\pm0{,}2$ & 10 & 0{,}98\\
(1,4) & $13{,}2\pm0{,}3$ & 13 & 1{,}02\\
(1,5) & $15{,}7$ & 16 & 0{,}98\\
(1,8) & $23{,}9\pm0{,}3$ & 25 & 0{,}96\\
(2,3) & $10{,}1$ & 11 & 0{,}92\\
(2,4) & $13{,}32\pm0{,}14$ (float64) & 14 & 0{,}95\\
(2,5) & $16{,}2$ & 17 & 0{,}95\\
(2,8) & $25{,}1\pm0{,}2$ (float64) & 26 & 0{,}97\\
(3,6) & $20{,}3\pm0{,}3$ & 21 & 0{,}97\\
(3,9) & $28{,}6$ & 30 & 0{,}95\\
(4,8) & $27{,}3\pm0{,}6$ & 28 & 0{,}98\\
\bottomrule
\end{tabular}
\end{center}
Alle Werte liegen innerhalb $\pm5\,\%$ an
\begin{equation}\label{eq:3dk}
\boxed{\;a_\sigma(k,d)\;=\;3d+k\;}
\end{equation}
--- unabhängig von $D=k(d-k)$: $(1,4)$ und $(2,4)$ liefern trotz
$D$-Unterschied um Faktor $3$ denselben Wert $13{,}2$.
\end{befund}

\subsection{Zur Frage nach dem Mechanismus der Quadratform}

Die Quadratform $\eta_{\mathrm{lo}}=a_\sigma\sigma^2$ mit
$a_\sigma=3d+k$ ist empirisch gesichert (Befund~\ref{bef:sigma},
\ref{bef:3dk}); ihre \emph{Herleitung} ist offen. Der naheliegende
AR(1)-Rate-Balance (linearisierte Form-Moden mit Relaxationsrate
$\gamma_{\min}\eta$ gegen Rauschdiffusion $\sigma^2$, stationäre Varianz
$\sigma^2/(2\eta\gamma_{\min})\ll1$, also
$a_\sigma=1/(2\gamma_{\min})$) ist durch direkte Messung der
Form-Moden-Dekadenz \textbf{widerlegt} worden:

\begin{befund}[Universelle lineare Kontraktionsrate]\label{bef:gamma}
Deterministische Dekadenz einer Form-Störung (mittelwert-subtrahierte
tangentiale Einheitsstörung) des exakt synchronisierten Zustands bei
$\sigma=0$, float64: Der frühe Anstieg von $\ln S_t$ ($S_t$ = chordale
Streuung) ist $-\gamma\eta$ mit
\[
\gamma(1,2) = 1{,}01,\ 1{,}01,\ 1{,}03
\quad\text{bei }\eta=0{,}01,\,0{,}02,\,0{,}05,
\qquad
\gamma(2,8) = 1{,}03,\ 1{,}05,\ 1{,}14
\quad\text{bei }\eta=0{,}05,\,0{,}1,\,0{,}2 .
\]
Die Rate ist also \emph{geometrieunabhängig} $\gamma\approx1$ ---
während die naive Heuristik
$\gamma_{\min}=1/(2a_\sigma)$, also $0{,}0195$ für $(2,8)$ und $0{,}071$
für $(1,2)$, verlangt hätte (Faktoren 50 bzw.\ 14 daneben, und
\emph{verschiedene} Raten für die beiden Geometrien).
\end{befund}

Der $\sigma^2$-Koeffizient wird folglich nicht durch die lineare
Relaxationsrate bestimmt, sondern durch die nichtlineare Geometrie des
rauschgetragenen Molekels gegenüber der Einfangsgrenze des geordneten
Beckens (Offenes Problem~\ref{prob:3dk}).

\begin{bemerkung}[Verhältnis zur Konjektur]
Die Konjektur postulierte $-a_\sigma\sigma^2$ (Verschiebung nach unten).
In Step-Size-Konvention schiebt Rauschen die untere Kante nach \emph{oben}
(+); eine Verschiebung nach unten entspricht entweder der oberen Kante
oder einer inversen $\eta$-Konvention. Die \emph{Form} (Quadrat in
$\sigma^2$, linear additiv) ist exakt bestätigt.
\end{bemerkung}

% =====================================================================
\section{Das obere Fensterende: Universalität für $k\ge2$}

\begin{befund}[Universelles $\eta_{\mathrm{hi}}$]\label{bef:hi}
Bei $\sigma=0$, Synchron-Start mit Störung der Norm $10^{-4}$, adiabatische
Aufwärtssweep:
\begin{center}
\begin{tabular}{lc}
\toprule
Konfiguration & $\eta_{\mathrm{hi}}$\\
\midrule
$(2,4)$, $(2,8)$, $(3,9)$, $(4,8)$ & $2{,}036$ (alle identisch bis Gitterauflösung $\pm0{,}03$)\\
\midrule
$(1,2)$, $(1,3)$, $(1,4)$, $(1,5)$, $(1,8)$ & $2{,}07$--$2{,}55$ (Median $2{,}23$; metastabil, siehe unten)\\
\bottomrule
\end{tabular}
\end{center}
Für $k\ge2$ ist $\eta_{\mathrm{hi}}$ also \emph{unabhängig von $(k,d)$}.
Für $k=1$ ist die Grenze ausgedehnt-metastabil: über 16 Seeds streut sie
zwischen $2{,}07$ und $2{,}55$, und $6/16$ Seeds bei $(1,2)$ desynchronisieren
im Bereich $\eta\le3{,}4$ überhaupt nicht.
\end{befund}

\begin{korollar}[Fenster-Schlussgesetz]\label{kor:schluss}
Die geordnete Phase existiert genau dann, wenn
\[
\eta_{\mathrm{lo}}(\sigma)\;<\;\eta_{\mathrm{hi}}
\quad\Longleftrightarrow\quad
\sigma\;<\;\sigma_{\max}:=\Big(\frac{\eta_{\mathrm{hi}}}{a_\sigma}\Big)^{1/2}.
\]
Für $(2,8)$: $\sigma_{\max}=(2{,}036/25{,}7)^{1/2}=0{,}281$;
beobachtet: geordnet bei $\sigma=0{,}2$, ungeordnet bei $\sigma=0{,}28$ --- konsistent.
Das Dichotomie-Fenster ist damit eine endliche Ressource, die das Rauschen
auf quantitativ vorhersagbare Weise zerstört.
\end{korollar}

% =====================================================================
\section{Anker: Mittlerer Drift, Stabilisierung, Widerlegung von $\lambda_W^2p_{\mathrm{fix}}$}

\begin{satz}[Mittlerer Ankerdrift]\label{thm:ankerdrift}
Für den Anker-Mix $\mu_p=p_{\mathrm{fix}}\,\delta_{W_{\mathrm{fix}}}+(1-p_{\mathrm{fix}})\,\mathrm{Haar}$ gilt
\[
\E_{W\sim\mu_p}\Big[(I-QQ^\top)\,W W^\top Q\Big]
\;=\;p_{\mathrm{fix}}\,(I-QQ^\top)\,W_{\mathrm{fix}}W_{\mathrm{fix}}^\top Q .
\]
\end{satz}
\begin{proof}
Für Haar-$W$ auf $\Gr(k,d)$ ist $\E[WW^\top]=\frac{k}{d}I_d$, und
$(I-QQ^\top)\frac{k}{d}I_d\,Q=\frac{k}{d}(Q-Q)=0$. Übrig bleibt nur der
Fix-Anteil.
\end{proof}

\begin{folgerung}
Jede driftgetriebene Verschiebung einer kritischen Kante ist in führender
Ordnung eine Funktion $F(\lambda_W p_{\mathrm{fix}})$ --- nicht von
$\lambda_W^2p_{\mathrm{fix}}$. Der Anker-Gemisch-Zweig (Haar) trägt zum
Mittel nicht bei, nur zur Varianz ($\propto\lambda_W^2p(1-p)$).
\end{folgerung}

\begin{befund}[Stabilisierung und Kantenverschiebung]\label{bef:anker}
Geordnete, verankerte NESS, $(2,8)$, $\sigma=0{,}05$: bei $\eta=0{,}3$
gilt $\lambda_{\mathrm{det}}=-0{,}021\,/\,-0{,}050\,/\,-0{,}116$ für
$(\lambda_W,p)=(0{,}05;\,0{,}5),\,(0{,}1;\,0{,}5),\,(0{,}2;\,0{,}5)$ ---
d.\,h.\ der Anker macht aus der marginalen Mannigfaltigkeit einen
kontrahierenden Attraktor mit Rate $\approx-2{,}3\,\lambda_W p_{\mathrm{fix}}$.

Die Desynchronisationsgrenze ($\lambda_{\mathrm{det}}=0$-Durchgang, scharf,
fehlerarm) steigt mit $\lambda_W p_{\mathrm{fix}}$ und sättigt:
\begin{center}
\begin{tabular}{lccc}
\toprule
$\lambda_W p_{\mathrm{fix}}$ & 0{,}025 & 0{,}05 & 0{,}1--0{,}2\\
Grenze & 1{,}77 & 2{,}03--2{,}09 & 2{,}11--2{,}15\\
\bottomrule
\end{tabular}
\end{center}
\textbf{Entscheidungstest (25$\sigma$):} konstruierte Paare mit
\emph{gleichem} $\lambda_W^2p_{\mathrm{fix}}=0{,}01$, aber verschiedenen
$\lambda_W p_{\mathrm{fix}}$ liefern Grenzen $2{,}094\,/\,2{,}116\,/\,2{,}146$ ---
untereinander unverträglich. Paare mit gleichem
$\lambda_W p_{\mathrm{fix}}=0{,}05$ liefern $2{,}085\,/\,2{,}094\,/\,2{,}111$
und alle fünf Testpunkte ordnen sich korrekt monoton in
$\lambda_W p_{\mathrm{fix}}$. Damit ist der konjizierte Term
$a_W\lambda_W^2p_{\mathrm{fix}}$ \textbf{widerlegt}; die Verschiebung ist
(gesättigt-)linear im mittleren Ankerdrift.
\end{befund}

% =====================================================================
\section{Marginal-Entartung: der Lyapunov-Beweis}

\begin{definition}[Deterministischer Ein-Schritt-Exponent]\label{def:ladet}
Sei $F_{\mathrm{det}}$ der deterministische Teil von \eqref{eq:dyn}
(ohne $G$; Anker-Ziehungen als Teil der Trajektorie). Für Einheitsblöcke
$U_i\in T_{Q_i}$, $i=1,\dots,N$, definiere
\[
\lambda_{\mathrm{det}}
\;=\;\lim_{T\to\infty}\frac1T\sum_{t<T}\;\frac1N\sum_{i=1}^N
\log\frac{\big\|F_{\mathrm{det}}(Q_i+\delta U_i)-F_{\mathrm{det}}(Q_i)\big\|}{\delta},
\qquad \delta\to0 .
\]
(Numerisch: Vorwärtsdifferenz, $\delta=10^{-9}$, float64; das Rauschen
treibt nur die Referenztrajektorie, nicht den Störungskanal.)
\end{definition}

\begin{befund}[Marginalität für $d\ge4$, Kontraktion für $d=2$]\label{bef:marginal}
$\sigma=0$, $\lambda_W=0$, exakt synchronisierter Zustand, float64:
\begin{itemize}
\item[$d\ge4$] ($(1,4),(2,4),(2,8)$):
\[
\lambda_{\mathrm{det}}(\eta)\;=\;0\;\pm\;0{,}0007
\quad\text{für \emph{alle}}\ \eta\in[0{,}002,\ 1{,}8],
\]
dann \emph{abrupter} Sprung: $-0{,}18$ bei $\eta=1{,}83$, $+0{,}405$ bei
$\eta=2{,}5$. Es gibt keinen weichen Eigenwert-Durchgang: das Plateau
endet diskontinuierlich.
\item[$d=2$] ($(1,2)$):
$\lambda_{\mathrm{det}}(\eta)<0$ strikt und wachsend:
$-0{,}002$ bei $\eta=0{,}002$, $-0{,}045$ bei $0{,}044$,
$-1{,}29$ bei $0{,}72$; Desynchronisation bei $2{,}2\pm0{,}15$.
\end{itemize}
Die Einteilung hängt an $d$, nicht an $k$: $(1,4)$ ist marginal wie
$(2,4)$; nur $(1,2)$ kontrahiert.
\end{befund}

\begin{bemerkung}[Bedeutung]
Auf der $\lambda_Wp_{\mathrm{fix}}=0$-Schicht ist die geordnete Phase für
$d\ge4$ \emph{exakt marginal} --- die top-Transversalrichtung weder zieht
an noch ab. Das ist die operationale Gestalt der behaupteten
Marginal-Entartung $h_1\equiv0$: der lineare Term der Bifurkations-Normalform
verschwindet identisch auf dem geordneten Blatt. Für $d=2$ ist der lineare
Term strikt kontrahierend (kontinuierlicher Charakter). Ankern hebt die
Marginalität auf (Satz~\ref{thm:ankerdrift}, Befund~\ref{bef:anker}):
erst mit Ankern existiert eine scharfe Grenze als
$\lambda_{\mathrm{det}}=0$-Durchgang.
\end{bemerkung}

% =====================================================================
\section{Die Hoeffding-Signatur $\zeta_1^\ast$}

\begin{befund}[$\zeta_1$ über die Phasen]\label{bef:zeta}
Mit dem Schätzer \eqref{eq:zetaest} (Kern $h=d_G$):
\begin{enumerate}
\item \textbf{Reine Phasen:} $\zeta_1=0$ tief in der geordneten Phase und
in der ungeordneten NESS (nach Subtraktion des Stichproben-Bodens
$\sigma_h^2/(N-1)$).
\item \textbf{Am Rand:} $\widehat\zeta_1$ springt auf einen universellen Wert:
\[
\zeta_1^\ast = 0{,}230\pm0{,}011\ \ ((2,8),\ \text{7 }\sigma\text{-Werte}),
\qquad
\zeta_1^\ast = 0{,}229\pm0{,}017\ \ ((1,2),\ \text{9 }\sigma\text{-Werte}).
\]
$N$-Skalierung am Rand: $0{,}217\,(N{=}32),\ 0{,}239\,(64),\ 0{,}246\,(128),\ 0{,}252\,(256)$
--- verträglich mit einer universellen Konstante $0{,}23$--$0{,}25$
(Trend nach oben: genuine $O(1/N)$-Inhomogenität jenseits des Bodens).
\item \textbf{Koexistenzbereich:} $\zeta_1\approx0{,}2$--$0{,}3$ durchgängig.
\item \textbf{Anker:} auch tief in der geordneten Phase
$\widehat\zeta_1\approx0{,}23\neq0$ (verankerte Wolke ist inhomogen).
Bedingung auf den \emph{momentanen} Anker-Modus (Fix vs.\ Haar) entfernt
$\approx0\,\%$; Regression von $r_i$ auf die persistente Wolken-Koordinate
(geodätischer Abstand zu $W_{\mathrm{fix}}$) entfernt $\approx74\,\%$.
Die bedingte Entartung $\zeta_1=0$ gilt also \emph{nur} auf der Schicht
$\lambda_Wp_{\mathrm{fix}}=0$.
\end{enumerate}
\end{befund}

Interpretation: $\zeta_1^\ast\approx1/4$ bedeutet, dass am kritischen Punkt
die erste Hoeffding-Komponente rund $23$--$25\,\%$ der Kernvarianz trägt ---
unabhängig von Geometrie und Rauschen. Ein analytischer Wert (Vermutung:
$1/4$) wäre ein scharfes Testkriterium für eine Theorie der Kante.

% =====================================================================
\section{Übergangscharakter}

\begin{befund}
\begin{enumerate}
\item \textbf{NESS bei $\sigma>0$:} kein auflösbarer Hysterese an keiner
Kante, weder $d=2$ noch $d\ge4$ (Zwei-Init-Protokoll: geordnete und
ungeordnete Startbedingungen liefern innerhalb der Fehler dasselbe
$\Delta(\eta)$; Auf/ab-Sweeps stimmen auf 1--3\,\%).
\item \textbf{$\sigma\to0$ (deterministisch):} der synchronisierte Zustand
ist deterministisch stabil weit unterhalb der ungeordnetseitigen Kante
(deterministisches Spinodal) --- es gibt eine Bistabilitätsstruktur, aber
die Rausch-Aktivierung zerstört sie bei $\sigma\ge0{,}05$ in
$<1{,}5\cdot10^4$ Schritten.
\end{enumerate}
Die klassische Klassifikation ,,kontinuierlich für $d=2$ / subkritisch
für $d\ge4$'' ist damit eine \emph{Null-Rauschen-Aussage}; im
rauschgetragenen NESS erscheinen beide Übergänge als kontinuierliche
Crossovers auf einer marginalen (nicht hysterese-getrennten) geordneten
Mannigfaltigkeit.
\end{befund}

% =====================================================================
\section{Bilanz gegenüber der ursprünglichen Konjektur}

\begin{center}
\begin{tabular}{p{6.1cm}lp{6.3cm}}
\toprule
Konjunkturelement & Urteil & Evidenz\\
\midrule
Quadratische $\sigma^2$-Verschiebung & \textbf{bestätigt} (Form) &
$\eta_{\mathrm{lo}}=a_\sigma\sigma^2$, Abschnitt $\approx0$, 12 Konfig.\\
Vorzeichen $-a_\sigma\sigma^2$ & konventionsabhängig &
Step-Size-Konv.: Rauschen hebt $\eta_{\mathrm{lo}}$; unten vs.\ oben\\
Skalierung von $a_\sigma$ & \textbf{widerlegt; ersetzt} &
gemessen $a_\sigma=3d+k$, unabhängig von $D\,d(d{+}2)/2$\\
Anker-Term $\lambda_W^2p_{\mathrm{fix}}$ & \textbf{widerlegt; ersetzt} &
$\lambda_Wp_{\mathrm{fix}}$ (mittl.\ Drift), sättigend; 25$\sigma$-Test\\
Marginal-Entartung auf $\sigma$-Blatt & \textbf{bestätigt} (operationell) &
float64-Lyapunov-Plateau $d\ge4$\\
Entartung unter Ankern erhalten & \textbf{widerlegt} &
Anker: $\lambda_{\mathrm{det}}<0$, $\zeta_1\approx0{,}23>0$\\
$\zeta_1=0$ bedingt & teilweise &
nur für $\lambda_Wp_{\mathrm{fix}}=0$; Regression entfernt 74\,\%\\
Kontinuierlich $d=2$ / subkritisch $d\ge4$ & bestätigt bei $\sigma\to0$ &
Kontraktion vs.\ Plateau; NESS-Hysterese fehlt bei jedem $\sigma>0$\\
\midrule
\multicolumn{3}{l}{\emph{Neu:} $\zeta_1^\ast\approx0{,}23$--$0{,}25$ universell;
Fenster-Schlussgesetz $\sigma_{\max}=(\eta_{\mathrm{hi}}/a_\sigma)^{1/2}$;
$\eta_{\mathrm{hi}}\approx2{,}04$ universell ($k\ge2$).}\\
\bottomrule
\end{tabular}
\end{center}

% =====================================================================
\section{Offene Probleme}\label{sec:offen}

\begin{vermutung}[Ursprung von $3d+k$]\label{prob:3dk}
Die lineare Kontraktionsrate ist universal $\gamma\approx1$
(Befund~\ref{bef:gamma}); die AR(1)-Heuristik scheidet damit aus.
Erkläre den Koeffizienten $a_\sigma=3d+k$ aus der nichtlinearen
Gleichgewichtsgeometrie: Kriterium vom Typ ,,Wolkenausdehnung der
ungeordneten NESS $\le$ Einfangsradius des geordneten Beckens'', mit
beiden Seiten als Funktion von $(\sigma;k,d,N)$. Erkläre insbesondere,
warum das Ergebnis additiv in $d$ und $k$ und unabhängig von
$D=k(d-k)$ ist. Gültigkeitsbereich: $\pm5\,\%$ für $d\le16$ bei
$N\ge64$; bei $d\ge24$ wächst $a_\sigma$ mit $N$
($\Delta a_\sigma\propto d\,\Delta\ln N$ bei $d=32$) --- die Grenze
dieser Gültigkeitsregion ist ebenfalls offen.
\end{vermutung}

\begin{vermutung}[Beweis von $\zeta_1^\ast=1/4$]
Berechne $\zeta_1$ am Rand exakt (Ein-Schleife um die ungeordnete NESS).
Gemessen: $0{,}2531$ ($N=512$), $0{,}2558$ ($N=1024$) --- konsistent mit
exakt $1/4$ im Limes $N\to\infty$; der $O(1/N)$-Abfall von $0{,}25$ bei
kleinem $N$ (gemessen $0{,}217$ bei $N=32$) ist mitzurechnen.
\end{vermutung}

\begin{vermutung}[Die Anker-Grenzfunktion]
Die Desynchronisationsgrenze $B(x)$, $x=\lambda_Wp_{\mathrm{fix}}$, ist
\emph{nicht-monoton}: kein Durchgang für $x\lesssim0{,}015$ (Marginalität
überlebt), steiler Anstieg, Maximum $B_{\max}\approx2{,}10$ bei
$x^\ast\approx0{,}075$, linearer Abfall mit Steigung $\approx-1{,}1$.
Herleite Schwellwert, Lage des Maximums und Abfallssteigung aus der
gestörten Plateau-Bedingung
$\lambda_{\mathrm{det}}=-\kappa\,\lambda_Wp_{\mathrm{fix}}+\ldots$
(gemessen $\kappa\approx2{,}3$) plus einem Verstimmungsterm
$\propto(\lambda_Wp_{\mathrm{fix}})^2$ für den Abfallzweig.
\end{vermutung}

\begin{vermutung}[Mechanismus der $d$-Abhängigkeit]
Warum ist das Plateau für $d\ge4$ exakt flach, für $d=2$ aber kontrahierend?
Kandidat: für $d=2$, $k=1$ ist der Tangentialraum eindimensional und alle
Relativ-Moden koppeln an den positive Spektrum eines Graph-Laplace-Typs;
für $d\ge4$ existieren isotrope Zusatzrichtungen, deren oberster Modus am
Nullpunkt hängt (Goldstone-artige Flachheit).
\end{vermutung}

\begin{vermutung}[Konventionsabbildung]
Weder $\eta_{\mathrm{lo}}$ noch $\eta_{\mathrm{hi}}$ bildet auf
$D\,d(d{+}2)/2$ unter einer $(k,d)$-unabhängigen Konstanten ab
($\eta_{\mathrm{hi}}$ ist für $k\ge2$ geometrieunabhängig; $a_\sigma=3d+k$).
Die $\eta$-Normierung der Quell-Theorie muss daher eine
$(k,d)$-abhängige Faktor enthalten; deren Bestimmung erfordert die dortige
Definition von $\eta$.
\end{vermutung}

% =====================================================================
\section{Analytische Herleitung der gemessenen Gesetze}\label{sec:herleitung}

Die fünf in Abschnitt~\ref{sec:offen} formulierten Probleme lassen sich
vollständig aus zwei exakten Strukturidentitäten der diskreten
Update-Regel~\eqref{eq:dyn} herleiten (ausführliche Tabellen und
Verifikationsdaten in \texttt{HERLEITUNG.md} und \texttt{verify\_theory.py}).

\begin{satz}[Exaktes Boltzmann-Gleichgewicht und Graßmann-Dualität]\label{thm:boltzmann}
Für $\lambda_W=0$ ist die Retraktion $X\mapsto\mathrm{orth}(Q+X)$ eine
Retraktion zweiter Ordnung auf $\Gr(k,d)$,
\[
P'=(Q+X)(I_k+X^\top X)^{-1}(Q+X)^\top
= P + (QX^\top+XQ^\top) + (XX^\top - QX^\top X Q^\top) + O(\|X\|^3),
\]
und der Drift $(I-P_i)\bar S\,Q_i=\mathrm{grad}_{Q_i}(-H)$ ist der exakte
Riemann'sche Gradient der Energie $H(P_1,\dots,P_N)=-\frac{1}{4N}\sum_{i,j}\mathrm{tr}(P_i P_j)$.
Im Kleinschritt-Limes $\eta,\sigma^2\to0$ bei festem $a:=\eta/\sigma^2$ ist die
stationäre Verteilung der Kette~\eqref{eq:dyn} das exakte Gibbs--Boltzmann-Maß
\begin{equation}\label{eq:boltzmann}
\rho_N(P_1,\dots,P_N)\;=\;\frac{1}{Z_N(a)}\,\exp\!\left(\frac{a}{2N}\sum_{i,j=1}^N \mathrm{tr}(P_i P_j)\right).
\end{equation}
Insbesondere hängen alle Gleichgewichtskanten nur von $a=\eta/\sigma^2$ ab
(F1) und erfüllen unter der Isometrie $P_i\mapsto I_d-P_i$ die exakte
Graßmann-Dualität $a_\sigma(k,d)=a_\sigma(d-k,d)$.
\end{satz}
\begin{proof}
Ein Vergleich mit der Riemann'schen Geodäten $\gamma(t)=QV\cos(t\Sigma)V^\top+U\sin(t\Sigma)V^\top$
($X=U\Sigma V^\top$) zeigt Identität von $P'$ und $\gamma(t)\gamma(t)^\top$ bis
$O(\|X\|^2)$. Der Fokker--Planck-Generator lautet daher
$\mathcal L f = \eta\langle\mathrm{grad}\,f,\mathrm{grad}(-H)\rangle+\frac{\sigma^2}{2}\Delta_{\Gr^N}f+O(\eta^2+\sigma^4)$;
partielle Integration gegen das Haar-Maß liefert $\mathcal L^\ast\rho_N=0$ für
$\rho_N\propto\exp(\frac{2\eta}{\sigma^2}(-H))$, d.\,h.\ \eqref{eq:boltzmann}.
\end{proof}

\begin{satz}[Lineare Weingarten-Instabilität --- Lösung von Problem~5]\label{thm:weingarten}
Für spurloses symmetrisches $M=\delta\bar S$ gilt unter dem Haar-Maß auf $\Gr(k,d)$:
\[
\E[\Delta P_{\mathrm{noise}}\mid Q]=-\sigma^2 d\,(P-\tfrac{k}{d}I_d),
\qquad
\E_{\mathrm{Haar}}[PM+MP-2PMP]=\frac{2k(d-k)}{(d-1)(d+2)}\,M.
\]
Die lineare Instabilitätsschwelle (oberes Spinodal) des ungeordneten Zustands ist daher
\begin{equation}\label{eq:alin}
a_{\mathrm{lin}}(k,d)\;=\;\frac{\eta_c^{\mathrm{lin}}}{\sigma^2}\;=\;\frac{d(d-1)(d+2)}{2D},
\qquad D=k(d-k).
\end{equation}
Für $k=1$ ($D=d-1$) ergibt sich exakt $a_{\mathrm{lin}}(1,d)=d(d+2)/2$.
Die Quell-Konstante $D\,d(d+2)/2$ bezeichnet somit das lineare Spinodal
$a_{\mathrm{lin}}$ (skaliert um $g(k,d)=D^2/(d-1)$ auf $\eta_{\mathrm{step}}/\sigma^2$),
nicht die Gleichgewichtskante $\eta_{\mathrm{lo}}$ oder die Flip-Kante $\eta_{\mathrm{hi}}$.
\end{satz}

\begin{satz}[Mittelfeld-Selbstkonsistenz, Ursprung von $3d+k$ und F8 --- Lösung von Problem~4]\label{thm:3dk}
Für den uniaxialen Ansatz $S(q)=\frac{k}{d}I_d+q(P_0-\frac{k}{d}I_d)$ ($h=aq$)
lautet die Mittelfeld-Freie-Energiedifferenz pro Agent
\[
\Delta F_a(q)\;=\;\frac{aD}{2d}\,q^2+\frac{k^2}{d}\,aq-\ln\E_{\mathrm{Haar}}\!\big[\exp(aq\,\mathrm{tr}(PP_0))\big].
\]
Im geordneten Bereich ($h\gg1$) liefert $\E_h[k-\mathrm{tr}(PP_0)]\simeq\frac{D}{2h}$
nach exakt herausfallender Tangentialdimension $D=k(d-k)$ die universelle
Zustandsgleichung
\begin{equation}\label{eq:selfcons}
q(1-q)\;=\;\frac{d}{2a}
\quad\Longrightarrow\quad
a_{\mathrm{sp}}=2d,\qquad
a_{\mathrm{1st}}(1,d)\underset{d\to\infty}{\sim}2{,}4554\,d.
\end{equation}
Die operationell gemessene Halb-Plateau-Kante $\Delta(\eta)=\frac12\Delta_{\mathrm{dis}}(k,d)$
ergibt parameterfrei $a_{1/2}^{\mathrm{eq}}(k,d)$, das alle 12 Konfigurationen
auf 1--4\,\% trifft und für $k\le d/2$ durch $3d+k$ (dualitätssymmetrisch:
$3d+\min(k,d-k)$) genähert wird. Der Bruch bei $d\ge16$ (F8) folgt aus der
thermischen Kramers-Nukleation $N\,\Delta F_{\mathrm{bar}}(a;k,d)=C(d,T)$ über
die Barriere des Übergangs erster Ordnung ($d\ge3$), was direkt
$\Delta a_\sigma\propto d\,\Delta\ln N$ und die Nicht-Vertauschbarkeit der
Limiten $N,d\to\infty$ liefert.
\end{satz}

\begin{satz}[Hoeffding-Signatur $\zeta_1^\ast\to1/4$ und $O(N^{-1})$-Bias --- Lösung von Problem~1]\label{thm:zeta1}
In jeder rauschgetragenen geordneten Phase ($\sigma>0$) bilden die Agenten eine
isotrope Tangentialwolke $X,Y\stackrel{\mathrm{iid}}{\sim}\mathcal N(0,\tau^2 I_D)$
mit $\|X-Y\|_2=\tau\sqrt{2D}+\tau(\frac12 u+\frac12 v-\frac{1}{\sqrt2}w)+O(D^{-1/2})$,
woraus unmittelbar
\[
\zeta_1(D)=\frac{\mathrm{Var}(g_1)}{\sigma_h^2}=\frac{1/4}{1/4+1/4+1/2}=\frac14-O(D^{-1})
\]
folgt (am Boltzmann-Rand für $(2,8)$: $\zeta_1^\infty=0{,}2588$). Der in
\eqref{eq:zetaest} definierte Schätzer besitzt unter der exakten
Hoeffding-Zerlegung $h=\theta+g_1(X)+g_1(Y)+g_2(X,Y)$ den endlichen-$N$-Bias
\begin{equation}\label{eq:zetabias}
\E[\widehat{\zeta_1}(N)]
\;=\;\zeta_1^\infty\left(1-\frac{5N-8}{N(N-1)}\right)-\frac{2}{N(N-1)},
\end{equation}
der die Messreihe $0{,}217\,(N{=}32)\to0{,}2558\,(N{=}1024)$ auf $\pm0{,}003$ reproduziert.
\end{satz}

\begin{satz}[Exakte Linearisierung, $d$-Dichotomie und Ankerkurve $B(x)$ --- Lösung von Problem~2 \& 3]\label{thm:lin}
Die exakte Ein-Schritt-Linearisierung von $F_{\mathrm{det}}$ am synchronisierten
Zustand $Q_1=\dots=Q_N=Q_0$ lautet für alle $(k,d)$ identisch
\begin{equation}\label{eq:linmap}
U_i'\;=\;(1-\eta)\,U_i+\eta\,\bar U
\qquad(\text{mit Anker } x=\lambda_W p_{\mathrm{fix}}:\ U_i'=(1-\eta R_{\mathrm{eff}}-x)\,U_i+\eta R_{\mathrm{eff}}\,\bar U),
\end{equation}
mit $D$ Goldstone-Eigenwerten $\mu_0=1$ ($\lambda_0=0$; mit Anker $\mu_0=1-x$)
und $(N-1)D$ Transversal-Eigenwerten $\mu_\perp=1-\eta$ (mit Anker
$\mu_\perp=1-\eta R_{\mathrm{eff}}-x$). Daraus folgen unmittelbar:
\begin{enumerate}
\item \emph{F3 \& F4:} Form-Moden-Rate $\frac{d\ln S_t}{dt}=\ln(1-\eta)$
(universell $\gamma=1$) und universelle Flip-Bifurkation bei $|1-\eta|=1\iff\eta_{\mathrm{hi}}=2$.
\item \emph{Problem~3 ($d$-Dichotomie):} Die Pro-Agent-Renormierung
$U_i\mapsto U_i'/\|U_i'\|_F$ auf $S^{D-1}$ friert bei $D=1$ ($(1,2)$,
$S^0=\{-1,+1\}$ unzusammenhängend) das Anfangs-$\bar U\approx0$ ein und misst
$\ln|1-\eta|$, während sie für $D\ge2$ alle $U_i$ im Bereich
$0<\eta<\eta_\ast(D):=\frac{2D}{D+1}$ ausrichtet ($\lambda_{\mathrm{det}}=0$)
und erst für $\eta>\eta_\ast(D)$ auf $\ln|1-\eta|$ abfällt ($\ln 1{,}5=+0{,}40547$
bei $\eta=2{,}5$).
\item \emph{Problem~2 (Ankerkurve $B(x)$):} $B(x)$ ist die untere Einhüllende aus
(i)~dem steigenden Ast $\varepsilon(\eta;\sigma)=x$ (mit $x_c\approx0{,}015$,
erzeugt durch fehlendes \texttt{gs\_orth} auf dem Output in \texttt{phaseE2.py})
und (ii)~der diskreten Flip-Grenze $\eta R_{\mathrm{eff}}+x=2\implies
B(x)=(2-x)/R_{\mathrm{eff}}$ mit Steigung $B'(x)=-1/R_{\mathrm{eff}}\approx-1{,}10$
(Schnittpunkt $x^\ast\approx0{,}075$, $B_{\max}\approx2{,}10$).
\end{enumerate}
\end{satz}

% =====================================================================
\section{Reproduktion}

Simulator: \texttt{diva.py} (GPU, PyTorch/CUDA; 3-d-BMM-Kerne,
Gram--Schmidt-Orthonormalisierung, elementweise Parameter, exakte
Hauptwinkel-SVD, Hoeffding-Maschine nach Korollar~\ref{kor:zeta},
deterministischer Differenzen-Exponent nach
Definition~\ref{def:ladet}). Messskripte \texttt{phaseA--I*.py},
\texttt{V1\_robustness.py} sowie \texttt{verify\_theory.py} (Verifikation aller
analytischen Herleitungen aus Abschnitt~\ref{sec:herleitung} und
\texttt{HERLEITUNG.md}); Rohdaten \texttt{results/*.json}; Figuren in
\texttt{results/figs/}. Gesamt-GPU-Aufwand: $\approx1$~h für das numerische
Programm plus $\approx1$~min für \texttt{verify\_theory.py} auf einer RTX~3060.

\end{document}
